\section{Introduction}

Stablecoins have emerged as a new instrument in parallel foreign exchange (FX) markets, particularly in economies with fixed or heavily managed exchange rates. In such settings, discrepancies between official and parallel exchange rates generate incentives for households and firms to seek alternative means of accessing foreign currency.

This paper develops a model in which stablecoins act as both (i) a lower-cost instrument for accessing FX and (ii) a source of public information about exchange rate misalignment. These two channels jointly affect equilibrium FX demand, crisis dynamics, and welfare.

The key insight is that stablecoins increase FX demand through both hedging and coordination motives, thereby increasing financial fragility even as they improve access to foreign currency. As a result, stablecoins can improve welfare at low levels of exchange rate misalignment but reduce welfare when misalignment is severe.
